\section*{Chapter \ref{ch:dv:no-arbitrage-and-the-fundamental-theorems} --- No Arbitrage and the Fundamental Theorems}

\begin{solution}{exo:dv:no-arbitrage-and-the-fundamental-theorems:1}
$q_u+q_d=0.99$ and $115q_u+95q_d=100$ give $20q_u=100-94.05$: $q_u=0.2975$,
$q_d=0.6925$. The risk-neutral up probability is $0.2975/0.99=0.3005$; the call
pays $(15,0)$ and is worth $15q_u=4.4625$.
\end{solution}

\begin{solution}{exo:dv:no-arbitrage-and-the-fundamental-theorems:2}
Below 5\,000: the calls pay nothing, the long 6\,000 put pays $6000-S$, the short
5\,000 put costs $5000-S$; total 1\,000. Between: the long call pays $S-5000$, the
long put $6000-S$; total 1\,000. Above 6\,000: the calls pay $(S-5000)-(S-6000)=1000$,
the puts nothing.
\end{solution}

\begin{solution}{exo:dv:no-arbitrage-and-the-fundamental-theorems:3}
The put pays $(0,0,20)$; its interval is $(20\times0,\ 20\times0.44)=(0,\ 8.80)$,
so 9.00 is an \omterm{def:dv:no-arbitrage-and-the-fundamental-theorems:arbitrage}{arbitrage}. Sell the put at 9.00 and super-replicate it with the
portfolio that matches it in the up and down states: $-0.5$ share and 60 bonds,
paying $(0,10,20)$ and costing $-50+58.8=8.80$. The edge is 0.20 today, plus 10 in
the middle state.
\end{solution}

\begin{solution}{exo:dv:no-arbitrage-and-the-fundamental-theorems:4}
A bond paying $(1,1)$ at 0.95 and an asset paying $(1,0)$ at price 0. The payoff
matrix is invertible, so each payoff has one \omterm{def:dv:no-arbitrage-and-the-fundamental-theorems:claim}{replicating portfolio} and one cost:
the \omterm{def:dv:no-arbitrage-and-the-fundamental-theorems:loop}{law of one price} holds. But buying the second asset costs nothing and pays 1
in the first state: an \omterm{def:dv:no-arbitrage-and-the-fundamental-theorems:arbitrage}{arbitrage}. The unique state-price vector is $(0,0.95)$, not
strictly positive.
\end{solution}

\begin{solution}{exo:dv:no-arbitrage-and-the-fundamental-theorems:5}
The risk-neutral probabilities are $q/0.98$, that is 0.3061, 0.4898 and 0.2041;
the \omterm{def:dv:no-arbitrage-and-the-fundamental-theorems:spd}{state-price density} is $q/\P$, that is 0.6667, 1.0667 and 2.0000. It is highest in the down state: a unit paid when
the share has fallen is worth twice its probability-weighted value, the price of
insurance against bad states.
\end{solution}

\begin{solution}{exo:dv:no-arbitrage-and-the-fundamental-theorems:6}
It is worth $q_u$, strictly between 0.10 and 0.54. At 0.54 it is super-replicated
by the portfolio paying $(1,0,0)$ in the up and down states: $\theta_s=1/40=0.025$
share and $-2$ bonds, which pays $(1,0.5,0)$ and costs $2.5-1.96=0.54$.
\end{solution}

\begin{solution}{exo:dv:no-arbitrage-and-the-fundamental-theorems:7}
Both vertices give 10.80: the interval is the single point 10.80. The spread pays
$(20,10,0)$, which is $0.5$ share minus 40 bonds in every state: it is replicable,
so its price is forced even though the market is incomplete.
\end{solution}

\begin{solution}{exo:dv:no-arbitrage-and-the-fundamental-theorems:8}
Inside the interval no price is an \omterm{def:dv:no-arbitrage-and-the-fundamental-theorems:arbitrage}{arbitrage} and none is ``fair'' without a
choice of \omterm{def:dv:no-arbitrage-and-the-fundamental-theorems:ad}{state prices}: 6 is the model's choice, not the market's. Bought at 5, the
call cannot be hedged to a sure profit, since the market cannot replicate it; the
best static hedge sub-replicates it for 2 and leaves the state-$u$ risk. The profit
of 1 exists only if the model's \omterm{def:dv:no-arbitrage-and-the-fundamental-theorems:ad}{state prices} are right, which is a bet.
\end{solution}

\begin{solution}{pb:dv:no-arbitrage-and-the-fundamental-theorems:1}
\textbf{1.} Long 5\,000 call $(S-5000)^+$, short 6\,000 call $-(S-6000)^+$, short
5\,000 put $-(5000-S)^+$, long 6\,000 put $(6000-S)^+$; the calls sum to
$\min(\max(S-5000,0),1000)$ and the puts to $1000-$ that amount: 1\,000.
\textbf{2.} USD~100\,000.
\textbf{3.} $973.40-324.60-156.70+470.00=962.10$ points, USD~96\,210.
\textbf{4.} $971.80-326.20-158.30+468.40=955.70$ points, USD~95\,570.
\textbf{5.} It removes all index risk; what remains is the counterparty (the
clearing house), the operational risk of a leg filled without the others, and
financing of the margin on a short box.
\textbf{6.} $-\ln(0.96210)=3.86\%$.
\textbf{7.} $-\ln(0.95570)=4.53\%$.
\textbf{8.} $958.90$; $-\ln(0.95890)=4.20\%$.
\textbf{9.} $1000-962.10=37.90$ points, USD~3\,790.
\textbf{10.} Lending $-13.6$~bp, borrowing $+53.1$~bp.
\textbf{11.} Legging pays half the width on each of four quotes (6.40 points in
all here); a package trades at one price near the mid, and removes the risk of a
partial fill.
\textbf{12.} The fair box is $1000e^{-0.04}=960.79$. Buying costs 962.10, more than
that, and selling brings 955.70, less: no \omterm{def:dv:no-arbitrage-and-the-fundamental-theorems:arbitrage}{arbitrage} (\texttt{check\_box} returns
nothing).
\textbf{13.} 960.79 points.
\textbf{14.} The short legs could be assigned early, turning the box into an open
position; the payoff would no longer be sure.
\textbf{15.} Treasuries carry a convenience yield (their use as collateral and
liquidity), about 40~bp on average in the study cited, so their yield is below the
riskless rate that boxes reveal.
\textbf{16.} Leveraged holders of index portfolios, funds and market makers who
want to borrow against margin at a rate close to the riskless one rather than at a
broker's margin rate.
\textbf{17.} $10\,000\,000/95\,890=104.3$: 104 boxes, paying USD~9\,972\,560 today and
receiving USD~10\,400\,000 in a year.
\textbf{18.} Positions cleared at the clearing house, which ports them to another
member or closes them out; her claim is on the clearing house, whose default
resources stand behind it.
\textbf{19.} 4.20\%, 19.7~bp above the 4.00\% overnight rate.
\textbf{20.} Because it pays the strike distance in every state, and a sure amount
at $T$ is exactly what a zero-coupon bond pays.
\end{solution}

\begin{solution}{iq:dv:no-arbitrage-and-the-fundamental-theorems:1}
Long call and short put at $K_1$, short call and long put at $K_2$: two synthetic
forwards that net to the sure amount $K_2-K_1$ at expiry. Its price is
$P(0,T)(K_2-K_1)$, so it quotes a riskless rate for the option's maturity with the
clearing house as counterparty; it needs European exercise.
\iqlookfor{parity twice, and the rate read off the price.}
\end{solution}

\begin{solution}{iq:dv:no-arbitrage-and-the-fundamental-theorems:2}
No \omterm{def:dv:no-arbitrage-and-the-fundamental-theorems:arbitrage}{arbitrage} holds if and only if there is a probability equivalent to the real one
under which discounted traded prices are martingales (in continuous time: no free
lunch with vanishing risk, and a sigma-martingale). Its probabilities are
normalised \omterm{def:dv:no-arbitrage-and-the-fundamental-theorems:ad}{state prices}: they encode prices, including risk premia, not beliefs
about how likely states are.
\iqlookfor{the equivalence, and ``pricing weights, not forecasts''.}
\end{solution}

\begin{solution}{iq:dv:no-arbitrage-and-the-fundamental-theorems:3}
With zero rate, $110q_u+95(1-q_u)=100$ gives $q_u=1/3$: the digital is worth 0.3333.
\iqlookfor{solving for the risk-neutral probability, not using a real-world one.}
\end{solution}

\begin{solution}{iq:dv:no-arbitrage-and-the-fundamental-theorems:4}
State-price vectors solve $D^\top q=p$; the solution is unique exactly when the
payoffs span every state, which is completeness. A trinomial market with a share
and a bond is incomplete: option prices are only bounded, and every positive
state-price vector gives an arbitrage-free price; choosing one is a model.
\iqlookfor{linear algebra behind the theorem, and bounds instead of prices.}
\end{solution}

\begin{solution}{iq:dv:no-arbitrage-and-the-fundamental-theorems:5}
At mids the butterfly costs $12.00-13.00+0.90=-0.10$, a convexity violation. At
executable prices it costs $12.05-2\times6.45+0.95=0.10$: no \omterm{def:dv:no-arbitrage-and-the-fundamental-theorems:arbitrage}{arbitrage}. The
mids are not tradable prices; the question is whether the quotes are.
\iqlookfor{checking at bid and ask, not at mid.}
\end{solution}

\begin{solution}{iq:dv:no-arbitrage-and-the-fundamental-theorems:6}
Keep the chain sorted by expiry and strike with forward and discount factor per
expiry; on each quote update recheck only the inequalities that involve the
changed strike (bounds, neighbours for monotonicity and slope, the triples for
convexity, the boxes through it): a constant number of checks. Use executable
prices, tolerances of one tick, and debounce stale quotes; log violations with
the portfolio; run a full linear-programme check periodically.
\iqlookfor{incremental checks local in strike, executable prices, and a periodic
global pass.}
\end{solution}
